\documentclass[10pt,a4paper]{amsart}
\usepackage[margin=19mm]{geometry}
\usepackage{amsmath,amssymb,mathtools}
\usepackage[scr=rsfs,cal=euler]{mathalfa}
\usepackage{xfrac}
\usepackage[unicode,hidelinks,bookmarks=false]{hyperref}
\newcommand{\ie}{i.e.\ }
\newcommand{\cf}{cf.\ }
\newcommand{\resp}{respectively\ }
\newcommand{\Subsection}{\subsection}
\providecommand{\nobreakdash}{\nobreak\hbox{-}}
\setlength{\emergencystretch}{2em}
\multlinegap0pt
\usepackage{bm}
\usepackage{bbm}
\usepackage{dsfont}
\usepackage{xspace}
\usepackage{cancel}
\usepackage{mathtools}
\usepackage{amsthm}
\makeatletter
\@ifundefined{Thm}{\newtheorem{Thm}{Thm}}{}
\@ifundefined{Lem}{\newtheorem{Lem}[Thm]{Lemma}}{}
\@ifundefined{Prop}{\newtheorem{Prop}[Thm]{Proposition}}{}
\makeatother
\title[An alternative approach for the mean-field behaviour of weak]{An alternative approach for the mean-field behaviour of weakly self-avoiding walks in dimensions $d>4$}
\author{Hugo Duminil-Copin, Romain Panis}
\date{Source: arXiv:2410.03649v2 (CC BY 4.0)}
\begin{document}
\maketitle

\noindent\emph{Source and licence.} An alternative approach for the mean-field behaviour of weakly self-avoiding walks in dimensions $d>4$. Version arXiv:2410.03649v2 (CC BY 4.0), distributed under the Creative Commons Attribution 4.0 International licence, as recorded in the arXiv licence metadata for this exact version and shown on the abstract page https://arxiv.org/abs/2410.03649v2. Full rights notice: licenses/arxiv-2410.03649-CC-BY-4.0.txt.\par\smallskip

\noindent\textit{Editorial scope.} Counted unit: the Lem whose source label is \texttt{lem:lower below full-space} in the article (source0.tex lines 960--965, source label \texttt{lem:lower below full-space}), delivered together with the article's own complete original proof of it (source0.tex lines 968--1027, one \texttt{proof} environment of 60 lines). No mathematics is written, repaired or re-derived: statement and proof are the article's own text line for line. Required context retained uncounted: eq:bound full plane (lines 118--122); eq:bound half plane (lines 118--122); Lem: Simon-Lieb WSAW (lines 211--221); eq: half plane at distance $k$ from half space (lines 302--306); eq: sum at distance k from halfplane (lines 302--306); prop: theorem beta start = beta\_c (lines 530--543); lem:lower below half-space (lines 916--921). Every definition, display and label of those lines is retained unchanged. Omitted from this package and not redistributed: 1--117, 123--210, 222--301, 307--529, 544--915, 922--959, 966--967, 1028--1383. Nothing inside a retained excerpt is omitted or altered: every retained body is delivered exactly as the article's lines are written, so no omission receipt of any kind accompanies this package. Citations: the retained excerpts carry no citation command at all, so no bibliography entry of the article is transcribed and no reference is redistributed. Reference adaptation: none was needed and none was made; every reference inside the delivered unit resolves inside it, so no unresolved reference or citation is produced. Printed slips kept literal: no printed word, symbol, hypothesis or number of the counted statement or of its proof was altered. Layout and class: the article's own class is \texttt{article} with packages including amsfonts, amsthm, amsmath, amssymb, amscd, mathrsfs, eucal, graphicx, epstopdf, pict2e, epic, xcolor, float, cite; the wrapper is the lane's pinned \texttt{amsart} editorial wrapper at 10pt on A4, which restates the theorem-like environments the retained text needs and adds the article's own macro definitions that the retained text needs (none), with extra wrapper packages (bm, bbm, dsfont, xspace, cancel, mathtools). The delivered numbering is the unit's own. Assets: no retained span contains a figure environment, a graphics-inclusion command, a PDF-page inclusion, a table or a TikZ picture; no image file is copied into this package, the package declares no asset dependency and no image file is redistributed with it. Licence: version arXiv:2410.03649v2 of this article is distributed under the Creative Commons Attribution 4.0 International licence, as recorded in the arXiv licence metadata for this exact version and shown on the abstract page https://arxiv.org/abs/2410.03649v2. A full rights notice is delivered at licenses/arxiv-2410.03649-CC-BY-4.0.txt. Characters: every retained excerpt is pure ASCII, so the delivered engine, XeLaTeX with its default Latin Modern text font, reports no missing character.
\begin{align}
\label{eq:bound full plane}
G_\beta(0,x)&\le \frac{C}{(1\vee |x|)^{d-2}}\exp(-|x|/L_\beta),\\
\label{eq:bound half plane} G_\beta^\mathbb H(0,x)&\le \frac{C}{(1\vee|x_1|)^{d-1}}\exp(-|x_1|/L_\beta).
\end{align}

\begin{Lem}\label{Lem: Simon-Lieb WSAW} Let $d\geq 2$. For $0<\beta<\beta_c$, $0\in S\subset \Lambda$ with $\Lambda\subset \mathbb Z^d$, and $x\in \Lambda$, 
\begin{align}\label{eq:SL}
G_\beta^\Lambda(0,x)&\le G_\beta^S(0,x)+ \sum_{\substack{y\in S\\ z\in \Lambda\setminus S\\ y\sim z}}G_\beta^S(0,y)\beta G_\beta^\Lambda(z,x),\\
G_\beta^\Lambda(0,x)&\ge G_\beta^S(0,x)+\sum_{\substack{y\in S\\ z\in \Lambda\setminus S\\ y\sim z}}G_\beta^S(0,y)\beta G_\beta^\Lambda(z,x)- \lambda \sum_{u\in S}E_\beta^{S,\Lambda}(u)G_\beta^\Lambda(u,x),\label{eq:reversed SL}
\end{align}
where 
\begin{equation}
E_\beta^{S,\Lambda}(u):=\sum_{\substack{y \in S\\ z\in \Lambda\setminus S\\ y\sim z}}G_\beta^S(0,u)G_\beta^S(u,y) \beta G_\beta^\Lambda(z,u).
\end{equation}
%$B(\beta):=\sum_{u\in \mathbb Z^d}G_\beta(u)^2$ and $\bfC(\beta):=4G_\beta(0)^4$. 
\end{Lem}

\begin{align}\label{eq: full plane estimate from half plane}
    G_\beta(0,x)&\le \frac{3C}{(1\vee|x|)^{d-2}}\qquad&\forall x\in \mathbb Z^d,\\ \label{eq: half plane at distance $k$ from half space}
    G_\beta^{\mathbb H_n}(0,x)&\le \frac{2C(k+1)}{(n-k)^{d-1}}\qquad&\forall x\in \partial\mathbb H_{n-k}\text{ with }1\le k<n,\\ \label{eq: sum at distance k from halfplane}
    \sum_{x\in \partial\mathbb H_{n-k}}G_\beta^{\mathbb H_n}(0,x)\beta&\le 4(k+1)\qquad&\text{ with }n,k\ge0.
\end{align}

\begin{Prop}\label{prop: theorem beta start = beta_c}
Fix $d>4$. There exist $C,\lambda_0>0$, and $K=K(C,d)>0$ such that for every $\lambda<\lambda_0$,
\begin{align}
G_{\beta_c}(0,x)&\le \frac{3C}{(1\vee|x|)^{d-2}}&\forall x\in \mathbb Z^d,\label{eq:upper bound below L}\\
G_{\beta_c}^{\mathbb H}(0,x)&\le \frac{C}{(1\vee|x_1|)^{d-1}}&\forall x\in \mathbb H
\label{eq:upper half-spacebound below L},\\
\label{eq: bound phi_n final prop}
\varphi_{\beta_c}(\mathbb H_n)&\le 1+K\lambda &\forall n\ge0,\\
\sup\{\varphi_{\beta_c}(B):B\in \mathcal B\}&\le 1+K\lambda, &\\
E^{\mathbb H_n,\mathbb Z^d}_{\beta_c}&\leq K &\forall n\geq 0,
\\
\sup\big\{E^{B,\mathbb Z^d}_{\beta_c}:B\in \mathcal B\big\}&\leq K. &\label{eq:general bound error last prop}
\end{align} 
\end{Prop}

\begin{Lem}\label{lem:lower below half-space}
Fix $d>4$. There exist $c,\varepsilon_0,\lambda_0>0$ such that for every $\lambda<\lambda_0$, every $\varepsilon<\varepsilon_0$, every $\tfrac1{2d}\le \beta\le \beta_c$, and every $x\in \mathbb H$ with $x_1=|x|\le 6L_\beta(\varepsilon)$,
\begin{equation}\label{eq:op}
G_\beta^{\mathbb H}(0,x)\ge \frac{c}{(1\vee|x|)^{d-1}}.
\end{equation}
\end{Lem}

\begin{Lem}\label{lem:lower below full-space}
Fix $d>4$. There exist $c,\varepsilon_0,\lambda_0>0$ such that for every $\lambda<\lambda_0$, every $\varepsilon<\varepsilon_0$, every $\tfrac1{2d}\le \beta\le \beta_c$, and every $|x|\le5L_\beta(\varepsilon)$,
\begin{equation}\label{eq:lb full below 5L}
G_\beta(0,x)\ge \frac{c}{(1\vee|x|)^{d-2}}.
\end{equation}
\end{Lem}

\begin{proof}Let $\lambda_0, \varepsilon_0$ be given by Lemma \ref{lem:lower below half-space}. We will choose $\lambda_0$ even smaller below. Let $\lambda<\lambda_0$ and $\varepsilon<\varepsilon_0$. By symmetry, we may consider $x\in A_n$, where $n=|x|$. Lemma~\ref{Lem: Simon-Lieb WSAW} applied to $S=\mathbb H_k$ and $\Lambda=\mathbb H_{k+1}$ gives that
\begin{multline}\label{eq: lower bound typical full space}
G_\beta(0,x)=G_\beta^{\mathbb H}(0,x)+\sum_{k\geq 0}\Big(G_\beta^{\mathbb{H}_{k+1}}(0,x)-G_\beta^{\mathbb H_k}(0,x)\Big)
\\\ge \sum_{k\ge 0}\sum_{\substack{y\in \mathbb H_k,\\ z\notin \mathbb H_k,\:y\sim z}}G_\beta^{\mathbb H_k}(0,y)\beta G_\beta^{\mathbb H_{k+1}}(z,x)-\lambda \sum_{k\geq 0}\sum_{u\in \mathbb H_k}E_\beta^{\mathbb H_k,\mathbb H_{k+1}}(u)G_\beta^{\mathbb H_{k+1}}(u,x).
\end{multline}
Looking at the first sum, we see that
\begin{align*}
 	\sum_{k\geq 0}\sum_{\substack{y\in \mathbb H_k,\\ z\notin \mathbb H_k,\:y\sim z}}G_\beta^{\mathbb H_k}(0,y)\beta G_\beta^{\mathbb H_{k+1}}(z,x)&\geq \sum_{k=0}^{n\wedge (L_\beta(\varepsilon)-1)}\frac1{2d}\varphi_\beta(\Lambda_k)\min\{G_\beta^{\mathbb H_{k+1}}(z,x):-z_1=|z|=k+1\}	
 	\\&\geq \sum_{k=0}^{n\wedge (L_\beta(\varepsilon)-1)}\frac1{2d}(1-\varepsilon)\frac{c}{(n+k)^{d-1}}
 	\\&\geq \frac{c_1}{n^{d-2}},
\end{align*}
where $c_1=c_1(d)>0$, and where we restricted our attention to special positions of $y$ and used both the bound $\varphi_\beta(\Lambda_k)\ge 1-\varepsilon$ provided by $k\le n \wedge (L_\beta(\varepsilon)-1)$ and the lower bound  from Lemma~\ref{lem:lower below half-space} (note that $n+n \wedge L_\beta(\varepsilon)\le 6L_\beta(\varepsilon)$ for $n\le 5L_\beta(\varepsilon)$).

Turning to the ``error'' term in \eqref{eq: lower bound typical full space}, we take the position of $u$ into account and use the upper bound \eqref{eq:bound half plane} to get 
\begin{align}
		\sum_{k\geq 0}\sum_{p\geq 0}\sum_{u\in \partial \mathbb H_{k-p}}\sum_{\substack{y\in \mathbb H_k\\z\notin \mathbb H_k\\y\sim z}} &G_\beta^{\mathbb H_k}(0,u)G_\beta^{\mathbb H_k}(u,y)\beta G_\beta^{\mathbb H_{k+1}}(z,u)G_\beta^{\mathbb H_{k+1}}(u,x)\nonumber
		\\&\leq \sum_{k\geq 0}\sum_{p\geq 0}\frac{C}{(p+1)^{d-1}}\cdot \varphi_\beta(\mathbb H_p)\cdot \sum_{u\in \partial \mathbb H_{k-p}}G_\beta^{\mathbb H_k}(0,u)G_\beta^{\mathbb H_{k+1}}(u,x).\label{eq: proof full space lower 1}
%		\\&\leq \sum_{k\geq 0}\sum_{p\geq 0}4(p+1)\cdot(1+\varepsilon)\cdot\frac{C}{(p+1)^{d-1}}\cdot \frac{C'(p\wedge(n+k))}{|n+k-p|^{d-1}+1}
%		\\&\leq \frac{C''}{n^{d-2}},
		\end{align}
Using Proposition \ref{prop: theorem beta start = beta_c}, we have that $\varphi_\beta(\mathbb H_p)\leq 1+K\lambda$. We bound \eqref{eq: proof full space lower 1} differently according to the value of $p$. 

First, use \eqref{eq:bound full plane} to get
\begin{align}
	\sum_{p\geq (n+k)/2}\frac{C(1+K\lambda)}{(p+1)^{d-1}}\sum_{u\in \partial \mathbb H_{k-p}}G_\beta^{\mathbb H_k}(0,u)G_\beta^{\mathbb H_{k+1}}(u,x)&\leq \frac{C_1}{(n+k)^{d-1}}\sum_{u\in \mathbb Z^d}G_{\beta}(0,u)G_{\beta}(u,x)
	\\&\leq \frac{C_2}{(n+k)^{d-1}},\label{eq: proof full space lower 2}
\end{align} where $C_1,C_2>0$ only depend on $C$ and $d$.
%where we used \eqref{eq: sum at distance k from halfplane}, Proposition \ref{prop: theorem beta start = beta_c}, and \eqref{eq: half plane at distance $k$ from half space}, and where $C',C''>0$ only depend on $C$ and $d$.
%The proof follows readily by choosing $\lambda_0$ small enough (so that $\lambda C''<c_1/2$).
%
%Restrict first to the case $p\geq n+k$ and use \eqref{eq: half plane at distance $k$ from half space} and \eqref{eq: sum at distance k from halfplane} to get that 
%\begin{equation}
%	\max_{u\in \partial \mathbb H_{k-p}}G_\beta^{\mathbb H_k}(0,u)\leq \frac{Ck}{(p-k)^{d-1}}, \qquad \sum_{u\in \partial \mathbb H_{k-p}}G_\beta^{\mathbb H_{k+1}}(u,x)\leq 4(n+k).
%\end{equation}
%Now, 
%\begin{equation}
%	\sum_{k\geq 0}\sum_{p\geq n+k}\frac{1}{(p+1)^{d-1}}\frac{k}{(p-k)^{d-1}}(n+k)=\sum_{p\geq n}\frac{1}{(p+1)^{d-1}}\sum_{k=0}^{p-n}\frac{(n+k)k}{(p-k)^{d-1}}\leq \sum_{p\geq n}\frac{1}{p^{d-3}}\sum_{k=0}^{p-n}\frac{1}{(p-k)^{d-1}} \leq \frac{C'}{n^{d-2}}.
%\end{equation}

Second, turning to the contribution coming from $p< \tfrac{n+k}{2}$, we write
\begin{equation}
	\max_{u\in \partial \mathbb H_{k-p}}G_\beta^{\mathbb H_{k+1}}(u,x)\stackrel{\eqref{eq: half plane at distance $k$ from half space}}\leq \frac{2^dC(p+2)}{(n+k)^{d-1}}, \qquad \sum_{u\in \partial \mathbb H_{k-p}}G_\beta^{\mathbb H_{k}}(0,u)\stackrel{\eqref{eq: sum at distance k from halfplane}}\leq \frac{4(p+1)}{\beta},
\end{equation}
where we used that for $u\in \partial \mathbb H_{k-p}$, one has $|u_1-x_1|=n+k-p\geq \tfrac{n+k}{2}$. Hence,
\begin{align}
	\sum_{p< (n+k)/2}\frac{C(1+K\lambda)}{(p+1)^{d-1}}\sum_{u\in \partial \mathbb H_{k-p}}G_\beta^{\mathbb H_k}(0,u)G_\beta^{\mathbb H_{k+1}}(u,x)&\leq C_3\sum_{p\leq\tfrac{n+k}{2}}\frac{1}{(p+1)^{d-3}}\frac{1}{(n+k)^{d-1}}
	\\&\leq \frac{C_4}{(n+k)^{d-1}},\label{eq: proof full space lower 3}
\end{align}
where $C_3,C_4>0$ only depend on $C$ and $d$. Combining \eqref{eq: proof full space lower 2} and \eqref{eq: proof full space lower 3} gives that
\begin{align}
	\sum_{k\geq 0}\sum_{p\geq 0}\sum_{u\in \partial \mathbb H_{k-p}}\sum_{\substack{y\in \mathbb H_k\\z\notin \mathbb H_k\\y\sim z}} G_\beta^{\mathbb H_k}(0,u)G_\beta^{\mathbb H_k}(u,y)\beta G_\beta^{\mathbb H_{k+1}}(z,u)G_\beta^{\mathbb H_{k+1}}(u,x)&\leq \sum_{k\geq 0}\frac{C_5}{(n+k)^{d-1}}\nonumber\\
	&\leq \frac{C_6}{n^{d-2}},
\end{align}
where $C_5,C_6$ only depend on $C$ and $d$. The proof follows by choosing $\lambda_0$ small enough so that $\lambda_0C_6<c_1/2$, and setting $c=c_1/2$.
%
%&\ge\sum_{k=0}^{\lfloor n/8\rfloor}\Big( \tfrac1{2d}\varphi_\beta(\Lambda_k)\min\{G_\beta^{\mathbb H_{k+1}}(z,x):-z_1=|z|=k+1\}-\frac{C_2\lambda}{n^{d-1}}\Big)\\
%&\ge \sum_{k=0}^{n}\Big(\tfrac1{2d}(1-\varepsilon)\frac{c}{(n+k)^{d-1}}-\frac{C_2\lambda}{n^{d-1}}\Big)\ge\frac{c'}{n^{d-2}}.
%
%In the second inequality, we used a bound on the error in this context, which resembles the same as the one obtained in the previous section, but which takes into account the position of $u$, as this position both impacts the bound on the error and on $G_\beta^{\mathbb H_{k+1}}(u,x)$ {\color{red}I will had some details this is not so easy from what I remember}. In the third, 
\end{proof}

\end{document}
